\section{问题一：典型日确定性储能调度}
\subsection{模型建立与求解}
给定全天负载、光伏和电价，以外网购电费为目标：
\begin{align}
\min\quad &J_1=\sum_{t=1}^{144}p_tg_t,\label{eq:q1obj}\\
\text{s.t.}\quad &g_t+V_t+b_t=L_t+c_t+w_t,\label{eq:q1balance}\\
&g_t\ge0,\qquad 0\le w_t\le V_t,\\
&E_0=E_{144}=6000,
\end{align}
并满足式\eqref{eq:storage}—\eqref{eq:bounds}。弃光变量使光伏富余时仍可保持能量平衡，首末状态相同则排除了消耗初始库存获得表面节省的情况。

模型首先按线性规划求解，再用二元变量$y_t$构造互斥对照：$c_t\le(5000/6)y_t$、$b_t\le(5000/6)(1-y_t)$。已有结果显示，本实例LP解的144个时段均不存在同时充放电，且LP与MILP目标相同。因此，对本实例而言，LP解已经是互斥模型的可行最优解。这里采用计算结果与约束核验作为依据，不把缺少边界条件的扰动论证推广到所有含弃光上界的模型。

\subsection{最优调度与费用}
\begin{table}[htbp]\centering\small
\caption{典型日储能调度与无储能基线}
\begin{tabular}{lrr}\toprule
方案 & 购电费/元 & 较基线节省/元\\\midrule
无储能 & 48052.05 & —\\
储能主方案 & 35126.95 & 12925.10\\
\bottomrule\end{tabular}\end{table}

无储能基线逐段优先利用光伏，不足部分由外网补齐，即$g_t^{\mathrm{base}}=\max(L_t-V_t,0)$。主方案购电费35126.95元，相对48052.05元基线下降26.90\%。调度过程中既有低价购电后储存，也有吸收光伏富余后在高价时段放电，经济作用来自全日约束下的联合安排。

\fig{01_q1_dispatch}
图\ref{fig:01_q1_dispatch}按同一时间轴对齐电价、供需、充放电与库存。储电量在允许区间内变化，日初和日末均为6000 kWh。低价时段充电与高价时段放电是总体特征，但局部行为还受光伏、负载、库存和功率限制影响，不能仅按电价阈值解释每个时段。

\subsection{费用变化的恒等分解}
为避免混淆效益归因，按逐段购电量变化定义
\begin{align}
A&=\sum_t p_t\max(g_t^{\mathrm{base}}-g_t,0),\\
B&=\sum_t p_t\max(g_t-g_t^{\mathrm{base}},0),\\
J_1&=J_1^{\mathrm{base}}-A+B.
\end{align}
少购电减少费用19111.64元，其他时段增购增加费用6186.54元，抵消后节省12925.10元。该分解直接由输入和调度复算，是现金费用的恒等关系，不声称把光伏消纳与套利的因果贡献完全分离。
\fig{02_q1_savings}
图\ref{fig:02_q1_savings}说明，部分时段增购电与全天节省并不矛盾：储能使电量从相对低成本时段流向高成本时段。完整逐段结果保留在正式交付表，正文给出题目指定的六个购电时段及四小时充放电汇总，见附录\ref{sec:outputs}。
\FloatBarrier
